% ================================================================
%  conteudo/introducao.tex
% ================================================================

\section{Introdução}

A crescente demanda por alternativas sustentáveis aos combustíveis fósseis
impulsiona o desenvolvimento de biocombustíveis de segunda geração, entre
os quais o combustível de aviação sustentável (\textit{Sustainable Aviation
Fuel}, SAF) ocupa posição estratégica. O setor aéreo é responsável por
aproximadamente 2,5\% das emissões globais de CO\textsubscript{2} e, ao
contrário de outros modais de transporte, dispõe de opções limitadas de
descarbonização de curto prazo, tornando o SAF uma das poucas rotas
tecnológicas viáveis para redução de emissões em escala (IATA, 2023).

As microalgas emergem como matéria-prima promissora para a produção de SAF
em razão de sua elevada produtividade volumétrica de lipídeos, capacidade de
crescimento em efluentes industriais e ausência de competição com culturas
alimentares por solo arável (CHISTI, 2007). Dentre as espécies estudadas,
\textit{Chlorella fusca} apresenta perfil lipídico rico em ácidos graxos de
cadeia média e longa -- C16 a C18 --, adequados à rota de conversão
\textit{Hydroprocessed Esters and Fatty Acids} (HEFA), tecnologia já
certificada pela norma ASTM~D7566 para blends de até 50\% em aviação
comercial (ASTM INTERNATIONAL, 2022).

O processo HEFA compreende duas etapas catalíticas principais: a
hidrodeoxigenação (HDO), que remove os grupos oxigenados dos triglicerídeos
via reação com hidrogênio, convertendo-os em \textit{n}-parafinas de cadeia
longa; e o hidrocraqueamento/isomerização, que ajusta a distribuição
carbônica e o ponto de congelamento do produto final às especificações de
combustível de aviação (PEARLSON et al., 2013). A qualidade do SAF obtido
depende diretamente do perfil de ácidos graxos da microalga, reforçando a
importância da caracterização lipídica da biomassa antes do processamento.

Um aspecto frequentemente negligenciado em estudos de viabilidade é o custo
do meio de cultivo, que pode representar parcela significativa do custo
operacional total. O aproveitamento de efluentes industriais ricos em
nitrogênio e fósforo -- como vinhaça diluída ou efluente de processamento de
alimentos -- como substrato parcial ou integral para o cultivo microalgal
constitui estratégia promissora de redução de custos e tratamento terciário
de resíduos simultaneamente (ABDEL-RAOUF et al., 2012).

Nesse contexto, o presente trabalho propõe o cultivo de \textit{Chlorella
fusca} em fotobiorreator por 15 dias, a caracterização do perfil lipídico da
biomassa produzida e a estimativa do rendimento de biojet fuel via simulação
do processo HEFA, acompanhada de análise orçamentária comparativa entre meio
sintético e meio parcialmente constituído por efluente industrial.

